\scanpage{56}
\mspage{V}
% archive foliation: 55
\uline{Cor} 1 Supposons que $G$ soit premier aux car. résiduelles.
\deleted{Alors \unsure{on}}
Supposons que pour tt sous-schéma affine intègre normal $S_0$ de $S$, et tt \uline{revêt} \uline{étale} $S'$ \inserted{intègre} de $S_0$, enfin tt sous-groupe $H$ de $G\times G$, on ait $H^1(S',H)$ fini.
Alors pour tt $X$ de t.f. sur $S$, $H^1(X,G)$ est fini.

Si $G$ est commutatif, il suffit de faire $H=G$ (\unsure{c'est} une condition \unsure{néce|nécessaire} et suffisante).

\uline{Corollaire} 2 Conditions équivalentes.

(i) Pour tt $X$ de t.f. sur $Y$, et tt groupe \inserted{fini} $G$ premier aux car. résiduelles de $X$, on a $H^1(X,G)$ fini.

(ii) \deleted{\unsure{Pour}} Idem avec $X$ revêt \deleted{fini} étale intègre d'un sous-schéma affine normal intègre $S_0$ de $S$.

\uline{NB} Dans le cor. 1, on a utilisé le fait que si $k$ est un corps alg. clos, alors pour tt $X$ de t.f. sur $k$ et $G$ premier à la car. de $k$, $H^1(X,G)$ est fini.
Voyons comment cela peut être prouvé par les méthodes envisagées ici.
